%=============================================================================!
% 1.1  QUANTITIES AND UNITS                                                   !
%=============================================================================!
\section{Quantities and units}
\label{sec:mo-units}

Every quantity in mechanics is a number multiplied by a unit. A height of
\SI{7.3}{\metre} is $7.3$ times the metre, and the same height expressed
in centimetres is \SI{730}{\centi\metre}; the number depends on the unit
chosen, whereas the height itself does not. This chapter requires only
two units, the metre for length and the second for time, and every other
unit it uses is built from these two, such as the metre per second
(\si{\metre\per\second}) for the rate at which distance is covered. In
this way of writing units, symbols set side by side are multiplied, and a
power $-1$ means one divided by that unit, so that $\mathrm{s}^{-1} =
1/\mathrm{s}$ and \si{\metre\per\second} is the metre divided by the
second.

The kind of quantity, as distinct from its size, is its dimension. We
write $\mathsf{L}$ for length and $\mathsf{T}$ for time, so that a rate of
covering distance, being a length divided by a time, has dimension
$\mathsf{L}\mathsf{T}^{-1}$. Since a length cannot meaningfully be added
to a time, only quantities of the same dimension can be added,
subtracted, equated or compared. The number that a function such as
$\sin$ acts on, called its argument, must be a pure number, one without
a unit, because the function is defined for numbers: the sine of three
metres has no meaning, and \cref{sec:mo-trig} shows that an angle
measured in radians is such a pure number. Every correct equation must
therefore balance in dimension, term by term. A car travelling at
\SI{20}{\metre\per\second} for \SI{3}{\second} covers $20\times 3 =
\SI{60}{\metre}$, and the dimensions of the product balance in the same
way as the numbers:
\begin{equation*}
   \underbrace{\mathsf{L}\mathsf{T}^{-1}}_{\text{rate}}
   \times\underbrace{\mathsf{T}}_{\text{time}} = \mathsf{L} .
\end{equation*}
This check costs nothing to apply. An equation that fails to balance is
certainly wrong, although one that balances is not thereby shown to be
right, since a missing pure number such as $\tfrac12$ leaves the
dimensions unchanged.

Atoms require much smaller units than everyday objects. Neighbouring atoms
are a few \aa ngstr\"oms apart (Chapter~8), where $\SI{1}{\angstrom} =
\SI{e-10}{\metre}$, and molecular dynamics advances in steps of about a
femtosecond, $\SI{1}{\femto\second} = \SI{e-15}{\second}$, for reasons
derived in Chapter~9. A speed of \SI{1}{\angstrom\per\femto\second} is
therefore
\begin{equation*}
   \frac{\SI{e-10}{\metre}}{\SI{e-15}{\second}}
   = 10^{-10-(-15)}\,\si{\metre\per\second}
   = \SI{e5}{\metre\per\second}.
\end{equation*}
